\documentclass[11pt]{article}
\usepackage[margin=1in]{geometry}
\usepackage[T1]{fontenc}
\usepackage[utf8]{inputenc}
\usepackage{lmodern,microtype,amsmath,amssymb,amsthm,mathtools}
\usepackage[colorlinks=true,linkcolor=blue,citecolor=blue,urlcolor=blue]{hyperref}
\hypersetup{pdftitle={Exact Gaussian energy and posterior identification of AND for ReLU networks},pdfauthor={ChatGPT-6 Astra; Kaarel Hanni; Dmitry Vaintrob}}
\newtheorem{theorem}{Theorem}
\newtheorem{lemma}[theorem]{Lemma}
\newtheorem{proposition}[theorem]{Proposition}
\newtheorem{corollary}[theorem]{Corollary}
\newtheorem{remark}[theorem]{Remark}
\newcommand{\R}{\mathbb R}
\newcommand{\E}{\mathbb E}
\newcommand{\Pp}{\mathbb P}
\newcommand{\relu}{\operatorname{ReLU}}
\newcommand{\norm}[1]{\lVert #1\rVert}
\title{Exact Gaussian energy and posterior identification\protect\\of AND for ReLU networks}
\author{ChatGPT-6 Astra \and Kaarel H\"anni \and Dmitry Vaintrob\thanks{Authors are listed in alphabetical order.}}
\date{7 November 2026}
\begin{document}
\maketitle

\begin{abstract}
We determine the exact minimum Euclidean variation norm of a shallow
ReLU score that represents AND on the full Boolean cube with signed
output margin at least $1/2$. For odd $q\ge3$ the value is
$(2q-1)/(2\sqrt{q+1})$; we give a separate closed formula for even $q$
and the two small dimensions. The corresponding minimum under the
isotropic Gaussian weight energy defined here is attained at every width
at least $1+\binom q2$ for $q\ge3$. A weighted two-design supplies this
polynomial width bound. Equality in an explicit dual certificate also
identifies every globally margin-valid energy minimizer after only the
positive all-plus input and its $q$ negative neighbors: all have the
AND sign function. Consequently, at each such fixed width, the actual
globally conditioned Gaussian posterior concentrates on AND in the
strong-prior limit after these $q+1$ labels. The target has a shallow
teacher of width $q+1$ with all weights at most one. These results do
not bound the number of online mistakes on arbitrary input orders.
\end{abstract}

\section{Model and results}

The global output-margin condition makes minimum-energy prediction a
problem over all unobserved Boolean labels as well as the observed
ones. A formula for the energy of a target alone therefore does not
identify the posterior after a training sample. We give both an exact
energy formula for AND and a sample on which every globally admissible
energy minimizer has the target's full sign function. The latter
statement concerns the whole posterior, including tied parameter
minimizers. It is not a mistake bound for arbitrary input orders.

For $q\ge1$, let $f_q:\{-1,1\}^q\to\{-1,1\}$ be positive only at
$x^0=(1,\ldots,1)$. Write $\bar x=(1,x)$ and consider
\begin{equation}\label{eq:network}
 F_{W,v}(x)=\sum_{j=1}^N v_j\relu(W_j\bar x),
 \qquad \relu(t)=\max\{t,0\}.
\end{equation}
The constant input is the only source of bias. There is one hidden
ReLU layer and a linear readout. The independent Gaussian weights are
\begin{equation}\label{eq:prior}
 W_{ji}\sim N(0,N\sigma^2/q),\qquad
 v_j\sim N(0,\sigma^2/N).
\end{equation}
Their density is proportional to $\exp[-K_N/(2\sigma^2)]$, where
\begin{equation}\label{eq:energy}
 K_N(W,v)=\frac qN\norm W_F^2+N\norm v_2^2.
\end{equation}
Before observing labels, condition on
\begin{equation}\label{eq:global}
 |F_{W,v}(x)|>\frac12\quad\text{for every }x\in\{-1,1\}^q.
\end{equation}
Then condition on exact observed signs. All strong-prior limits below
mean $\sigma\downarrow0$ at a fixed finite $q,N$.

For a real-valued function $F$ on the cube, define its Euclidean ReLU
variation norm by
\[
 \norm F_{\mathcal V}
 =\inf\left\{\sum_j |c_j|\norm{u_j}_2:
 F(x)=\sum_j c_j\relu(u_j\bar x)\text{ on the cube}\right\},
\]
where the infimum ranges over finite representations. Let
\[
 V_q=\inf\left\{\norm F_{\mathcal V}:
 f_q(x)F(x)\ge\frac12\text{ on the cube}\right\}.
\]

\begin{theorem}[Exact norm and Gaussian energy]\label{thm:energy}
The exact values are
\begin{equation}\label{eq:values}
 V_q=
 \begin{cases}
  1/\sqrt2,&q=1,\\[2pt]
  (\sqrt2+\sqrt6)/4,&q=2,\\[2pt]
  (2q-1)/(2\sqrt{q+1}),&q\ge3\text{ odd},\\[2pt]
  \dfrac{2q^2+\sqrt{q(q-2)}}{2\sqrt{q(q+1)(q+2)}},
    &q\ge4\text{ even}.
 \end{cases}
\end{equation}
Put $N_q=2$ for $q\le2$, and $N_q=1+\binom q2$ for $q\ge3$.
For every integer $N\ge N_q$, the minimum of \eqref{eq:energy} subject
to $f_q(x)F_{W,v}(x)\ge1/2$ on the full cube is exactly
\[
 2\sqrt q\,V_q.
\]
The strict-margin infimum has the same value. The norm lower bound
holds for every width; the stated width threshold is sufficient and
is not asserted necessary.
\end{theorem}

For $1\le i\le q$, let $x^i=x^0-2e_i$. Call the $q+1$ examples
$(x^0,+1),(x^1,-1),\ldots,(x^q,-1)$ the \emph{star sample}.

\begin{theorem}[Identification from the star sample]\label{thm:posterior}
Fix $N\ge N_q$. Every minimum-energy network satisfying the star
sample and the global weak absolute-margin condition has the full
sign function $f_q$. Under \eqref{eq:prior}--\eqref{eq:global}, after
conditioning on the star sample,
\[
 \Pp_\sigma\bigl(\operatorname{sign}F_{W,v}=f_q
       \text{ on the full cube}\mid\text{star sample}\bigr)
       \longrightarrow1.
\]
The probability remains at least this large after any further
target-consistent labels. In particular, at fixed $q,N$, the expected
number of subsequent sampling mistakes tends to zero on any
continuation without repetitions.
\end{theorem}

\begin{proposition}[A bounded teacher]\label{prop:teacher}
For $q\ge3$, $f_q$ has a depth-two teacher of width $q+1$ with every
individual weight at most one and strict global output margin
$q/[2(q-1)]>1/2$. For $q=1,2$, width two suffices with the same weight
bound. Thus the width thresholds in Theorems~\ref{thm:energy}
and~\ref{thm:posterior} are above these explicit bounded teacher
widths. A bounded-fan-in Boolean circuit for $f_q$ uses $q-1$ AND
gates when $q\ge2$.
\end{proposition}

\section{A dual certificate for all ReLU atoms}

We first treat $q\ge3$. Define
\begin{equation}\label{eq:parameters}
 \begin{gathered}
 n=q+1,\qquad k=\lceil n/2\rceil,\qquad \delta=2k-n\in\{0,1\},\\
 R=\sqrt{(q-1)^2-\delta^2},\qquad
 \varepsilon=\frac{q-1-R}{n},\qquad A=q-1-\varepsilon,\\
 D=2\sqrt{\frac{k(n-k)}n}=\sqrt{n-\frac{\delta^2}n}.
 \end{gathered}
\end{equation}
Here $R>0$, $A>0$, and $0\le\varepsilon<1$.
For a score $H$, put
\begin{equation}\label{eq:dual}
 T_A(H)=A H(x^0)-\sum_{i=1}^q H(x^i).
\end{equation}

\begin{lemma}[Atom bound and equality]\label{lem:atom}
For every $u=(b,w)\in\R^{q+1}$,
\begin{equation}\label{eq:atom}
 |T_A(\relu(u\bar x))|\le D\norm u_2.
\end{equation}
Among unit atom vectors, positive equality holds only at
\begin{equation}\label{eq:positive}
 H_+(x)=\frac1D\relu\!\left(-(1+\varepsilon)
                  +(1-\varepsilon)\sum_i x_i\right).
\end{equation}
Negative equality atoms all vanish at $x^0$ and have the form
\begin{equation}\label{eq:negative}
 H_S^{(j)}(x)=\relu\!\left(\frac2D
    \left[\frac jn\left(1+\sum_i x_i\right)
                    -\sum_{i\in S}x_i\right]\right),
 \quad |S|=j,
\end{equation}
where $j$ maximizes $j(n-j)$ among $0\le j\le q$.
Thus $j=k$ for odd $q$; for even $q$, either $j=q/2$ or $j=q/2+1$.
\end{lemma}
\begin{proof}
Write $s=b+\sum_iw_i$. If $s<0$, then $T_A\le0$. If $s\ge0$,
using $\relu(t)\ge t$ gives
\[
 T_A\le (A-q+2)s-2b
       =(-1-\varepsilon)b+(1-\varepsilon)\sum_iw_i.
\]
The coefficient vector $u_+$ in this last expression satisfies
\[
 \norm{u_+}_2^2
 =n-2(q-1)\varepsilon+n\varepsilon^2
 =n-\delta^2/n=D^2.
\]
Cauchy--Schwarz proves the positive bound.

For the negative bound, expand a sum of ReLUs as a maximum over
subsets $S\subseteq[q]$, and write $j=|S|$. Let
$e=(1,\ldots,1)\in\R^n$, and let $a_S\in\R^n$ be the indicator of
the selected input positions, zero at the constant position.
On the halfspace $s\le0$, the linear coefficient in $-T_A$ is
$je-2a_S$, whose inner product with $e$ is $j(q-1)\ge0$.
On $s\ge0$, the coefficient is $(j-A)e-2a_S$, and
\[
 \langle(j-A)e-2a_S,e\rangle
   =j(q-1)-An\le q(q-1)-An=-R<0.
\]
In either case, the component along $e$ contributes a nonpositive
inner product with the allowed atom vector. The perpendicular
projection of the coefficient is
\[
 -2(a_S-je/n),
 \qquad
 \norm{-2(a_S-je/n)}_2^2
     =\frac{4j(n-j)}n\le D^2.
\]
This proves the negative bound on every activation cone.

Positive equality forces equality in Cauchy--Schwarz, hence the
unit vector $u_+/D$. Its center preactivation is $R/D$, and its
neighbor preactivations are $[R-2(1-\varepsilon)]/D\ge0$.
The last inequality follows directly for $q=3,4$; for $q\ge5$ use
$R\ge q-2\ge3$. Thus it is indeed an equality atom.

For negative equality, the projection norm forces $j(n-j)$ to be
maximal. Such $j$ is nonzero, so both halfspace coefficients have
a strictly adverse component along $e$. Equality therefore requires
$s=0$, followed by equality in the perpendicular Cauchy--Schwarz
step. This gives exactly \eqref{eq:negative}. Conversely these
vectors have unit norm, vanish at $x^0$, and attain $-D$.
\end{proof}

Any score fitting the star margins satisfies
\begin{equation}\label{eq:lower}
 T_A(F)\ge\frac{A+q}{2},\qquad
 \norm F_{\mathcal V}\ge\frac{A+q}{2D}.
\end{equation}
This is a bound on every nonlinear competitor, not only a symmetric
or width-one representation.

\section{A full-cube witness with polynomially many atoms}

We use the negative atoms \eqref{eq:negative} with the upper size
$j=k=\lceil n/2\rceil$. Choose a distribution on these $k$-subsets
whose first and pair moments are
\begin{equation}\label{eq:moments}
 \Pp(i\in S)=\frac kq,\qquad
 \Pp(i,j\in S)=\frac{k(k-1)}{q(q-1)}\quad(i\ne j).
\end{equation}
The uniform distribution is one choice. Put $H_-=\E_S H_S^{(k)}$.

\begin{lemma}[The negative atom average]\label{lem:negative-average}
At every cube point other than $x^0$,
\[
 H_-(x)\ge D/q.
\]
Equality holds at each one-coordinate neighbor of $x^0$.
\end{lemma}
\begin{proof}
Suppose $h\ge1$ coordinates of $x$ are negative. Let $J$ count
how many of those positions belong to $S$, and put $Z=nJ-kh$.
Then
\[
 H_-(x)=\frac4{nD}\E Z_+.
\]
The random variable $Z$ is the sum of $h$ population entries,
each equal to $n-k$ at a selected coordinate and $-k$ otherwise.
Consequently $-kh\le Z\le(n-k)h$, so pointwise
\[
 Z^2\le h[(n-k)Z_++kZ_-]=h[nZ_+-kZ].
\]
The moments \eqref{eq:moments} give $\E Z=kh/q$ and
\begin{equation}\label{eq:moment-identity}
 \frac{\E Z^2}{nh}+\frac kn\E Z
 =\frac{k(n-k)}q
    +\frac{(h-1)k(2k-n)}{q(q-1)}
 \ge\frac{k(n-k)}q.
\end{equation}
For completeness, a population entry $Y_i=n1_{i\in S}-k$ has
\[
 \E Y_i^2=\frac{k[n^2-(q+2)k]}q,\qquad
 \E Y_iY_j=\frac{k[(q+3)k-n^2]}{q(q-1)}\quad(i\ne j),
\]
which yield \eqref{eq:moment-identity} by expanding the square of
the sum. Taking expectations in the pointwise bound proves
$\E Z_+\ge k(n-k)/q$. Since $D^2=4k(n-k)/n$, the claimed lower
bound follows. At $h=1$, the direct two-value calculation is an
equality.
\end{proof}

Set
\begin{equation}\label{eq:witness}
 \alpha=\frac D{2R},\qquad
 \beta=\frac qD\left(1-\frac{1-\varepsilon}R\right),\qquad
 F=\alpha H_+-\beta H_-.
\end{equation}
Both coefficients are positive. At the center, $H_-=0$ and
$H_+=R/D$, so $F(x^0)=1/2$. At any other point, $H_+$ is no larger
than its neighbor value. Lemma~\ref{lem:negative-average} therefore
gives
\[
 F(x)\le\frac\alpha D[R-2(1-\varepsilon)]-\frac{\beta D}q
       =-\frac12.
\]
The variation cost matches \eqref{eq:lower}:
\begin{equation}\label{eq:cost}
 \alpha+\beta=\frac{A+q}{2D}.
\end{equation}
One way to verify the identity is to use
$D^2-2q(1-\varepsilon)=-(1+\varepsilon)R$.
Substituting \eqref{eq:parameters} into \eqref{eq:cost} gives the
odd and even formulas in \eqref{eq:values}.

\begin{lemma}[Quadratic support size]\label{lem:design}
A distribution satisfying \eqref{eq:moments} can be supported on at
most $\binom q2$ subsets. Hence \eqref{eq:witness} can use at most
$1+\binom q2$ ReLU atoms.
\end{lemma}
\begin{proof}
Let $M=\binom q2$, and assign to each $k$-subset its pair-indicator
vector $v_S=(1_{i,j\in S})_{i<j}\in\R^M$. Every such vector lies in
the affine hyperplane $\sum_{i<j}(v_S)_{ij}=\binom k2$, of dimension
at most $M-1$. The uniform pair-moment vector is in their convex
hull, so Carath\'eodory's theorem represents it using at most $M$
vectors. Equivalently, start with the uniform distribution and
successively eliminate an affine dependence whenever its support
exceeds $M$.

Matching pair moments also matches first moments because $k\ge2$
and
\[
 \sum_{j\ne i}1_{i,j\in S}=(k-1)1_{i\in S}.
\]
This proves \eqref{eq:moments} with the claimed support.
\end{proof}

\begin{remark}[Hadamard improvement]
When $q$ is odd and a Hadamard matrix of order $q+1$ exists, normalize
its first row and column to be positive. Its $q$ remaining rows have
zero sum. Their negative input positions form $k=(q+1)/2$ subsets;
column orthogonality gives \eqref{eq:moments}. Thus $q$ negative
atoms and one positive atom suffice, reducing the width threshold
to $q+1$. No general Hadamard existence assumption is used in the
quadratic width bound.
\end{remark}

\section{Small dimensions}

For $q=2$, use $A=2/\sqrt3$ and $D=\sqrt{8/3}$ in
\eqref{eq:dual}. The negative halfspace-projection argument still
applies: the perpendicular norm is at most $D$, and
$j(q-1)-A(q+1)=j-2\sqrt3<0$ for $j\le2$.
For the positive bound, convexity gives
\[
 H(x^1)+H(x^2)\ge2\relu(b).
\]
When the center preactivation is nonnegative,
\[
 T_A(H)\le A(b+w_1+w_2)-2\relu(b)
          \le A(w_1+w_2)\le D\norm{(b,w)}_2,
\]
because $0<A<2$. A negative center preactivation gives $T_A\le0$.
Equality on the positive side holds only at
$H_+(x)=\relu(x_1+x_2)/\sqrt2$. Negative equality atoms vanish at
the center, by the same strict halfspace argument.

The two-atom witness
\[
 F(x)=\frac14\bigl[\relu(x_1+x_2)-\relu(2-x_1-x_2)\bigr]
\]
has outputs $1/2,-1/2,-1$ on the three Hamming layers. Its cost is
$(\sqrt2+\sqrt6)/4$, equal to the dual lower bound $(A+2)/(2D)$.

For $q=1$, the atom difference $H(1)-H(-1)$ is at most
$\sqrt2\norm{(b,w)}_2$ in absolute value. If both preactivations
are nonnegative, use $|b|\ge|w|$ to bound $2|w|$; if exactly one is
nonnegative, use Cauchy--Schwarz on $b\pm w$; if neither is
nonnegative, the difference vanishes. The two required margins
have difference at least one. The matching witness is
\[
 F(x)=\frac14\bigl[\relu(1+x)-\relu(1-x)\bigr]=x/2
 \quad\text{on }\{-1,1\}.
\]
Its variation cost is $1/\sqrt2$.

\section{Gaussian energy and posterior identification}

\begin{proof}[Proof of Theorem~\ref{thm:energy}]
The preceding dual bounds and witnesses prove all norm formulas.
For each neuron, arithmetic--geometric mean gives
\[
 \frac qN\norm{W_j}_2^2+Nv_j^2
       \ge2\sqrt q\,|v_j|\norm{W_j}_2.
\]
Summing proves the Gaussian lower bound. Conversely, for a unit
atom with coefficient $c\ne0$, choose a positive row scale $t$
with $t^2=N|c|/\sqrt q$, and use readout coefficient $c/t$.
Positive homogeneity preserves the atom's contribution, and its
energy is exactly $2\sqrt q\,|c|$. Apply this to the witnesses and
pad unused neurons by zeros. Lemma~\ref{lem:design} and the
two-atom small-dimensional witnesses give the stated thresholds.
Finally, readout amplification by factors tending down to one
proves equality of strict and weak infima.
\end{proof}

\begin{lemma}[Sign rigidity after the star sample]\label{lem:rigidity}
For $N\ge N_q$, every minimum-energy network fitting the star
sample and global weak absolute margins has sign function $f_q$.
\end{lemma}
\begin{proof}
For $q\ge3$, the full-cube witness and the star dual show that the
minimum energy is $2\sqrt q\,V_q$. Equality must hold throughout
\[
 K_N\ge2\sqrt q\sum_j|v_j|\norm{W_j}_2
       \ge\frac{2\sqrt q}D T_A(F)
       \ge2\sqrt q\,V_q.
\]
Every nonzero atom therefore attains dual equality, with its
coefficient sign agreeing with its dual sign. By
Lemma~\ref{lem:atom}, the score has the form
\[
 F=\alpha' H_+-\sum_S\beta_S H_S,
 \qquad \alpha',\beta_S\ge0,
\]
where each negative atom is nonnegative and vanishes at $x^0$.
Equality in the star-margin bound forces all star margins to equal
$1/2$, in particular $F(x^0)=1/2$. Thus $\alpha'=D/(2R)$.
At every other point, the strictly smaller value of $H_+$ gives
\[
 F(x)\le\alpha'H_+(x)<\alpha'H_+(x^0)=1/2.
\]
Global weak absolute margin forces $F(x)\le-1/2$.

For $q=2$, use the positive and negative equality cases from the
small-dimensional argument. The positive equality atom is zero
at every noncenter cube point, so the same conclusion holds.
For $q=1$, the star sample is the full domain.
\end{proof}

\begin{proof}[Proof of Theorem~\ref{thm:posterior}]
Let $\mathcal E$ be the event of strict global absolute margin
and the star labels. Its energy infimum is $E_*=2\sqrt q\,V_q$,
by the witnesses and readout amplification. Let $\mathcal B$ be
the closed set satisfying global weak absolute margins and the
star labels but assigning a positive label to some other cube
point. This is a finite union of closed sign cells. If nonempty,
it attains an energy minimum by coercivity. Lemma~\ref{lem:rigidity}
implies
\[
 \inf_{\mathcal B}K_N>E_*.
\]

For completeness, this strict energy gap controls the actual
finite-dimensional Gaussian posterior. Choose a strict feasible
point of energy below $\inf_{\mathcal B}K_N$, and a sufficiently
small ball around it contained in $\mathcal E$. Integration of
the Gaussian density over this ball gives a lower bound on
$\Pp_\sigma(\mathcal E)$ with a smaller exponential rate than the
Gaussian upper bound for the closed set $\mathcal B$. Thus
\[
 \Pp_\sigma(\mathcal B\mid\mathcal E)\longrightarrow0.
\]
The statement is trivial if $\mathcal B$ is empty. This proves
concentration on the full AND function without selecting any
particular parameter minimizer.

Additional target-consistent labels preserve the entire AND
function cell and can only remove other cells. Its posterior
probability therefore cannot decrease. A continuation without
repetitions has at most $2^q-q-1$ remaining inputs, so its expected
number of sampling mistakes is at most this number times the
post-star probability of a non-AND function, which tends to zero.
\end{proof}

\section{Teacher and scope}

\begin{proof}[Proof of Proposition~\ref{prop:teacher}]
For $q\ge3$, use $q$ identical hidden neurons
\[
 \relu\!\left(\frac{\sum_i x_i-(q-1)}{q-1}\right),
\]
each with readout coefficient one. Each is $1/(q-1)$ at the
center and zero elsewhere. Add a constant neuron $\relu(1)$
with readout coefficient $-q/[2(q-1)]$. The output is exactly
$q/[2(q-1)]$ at the center and its negative elsewhere. All
individual weights have magnitude at most one.

For $q=2$, the width-two score
\[
 \frac12\relu(x_1+x_2)-\relu\!\left(1-\frac{x_1+x_2}{2}\right)
\]
has outputs $1,-1,-2$ on the three Hamming layers and all weights
at most one. For $q=1$, use
$[\relu(1+x)-\relu(1-x)]/2=x$ on the cube.
The circuit statement is the usual binary AND tree.
\end{proof}

The norm values are intrinsic to ordinary ReLU and the isotropic
Gaussian energy defined here. Their noninteger form does not by
itself supply an ordering of exponentially many competing
functions. Nor does the displayed construction prove an additive
norm formula for multiple gadgets: an AND of disjoint AND blocks
is simply one AND. The posterior theorem applies once the star
labels have been observed. It leaves open how many mistakes can
occur before then on an arbitrary permutation of the cube.

\end{document}
